\section{Algoritmo}

Estas notas estudian algoritmos. Un \concept{algoritmo} es una secuencia de pasos. Pero nos interesan los algoritmos que resuelven un problema de computación, entonces:

\begin{definicion}{Algoritmo}{algoritmo}
Secuencia de pasos para transformar un conjunto de valores de \concept{entrada} (\textit{input}) a un conjunto de valores de \concept{salida} (\textit{output}).
\end{definicion}
No se escribe \say{a otro conjunto de valores} porque la salida podría ser igual a la entrada.

Un algoritmo se puede expresar en:
\begin{itemize}
  \item Lenguaje natural (español, inglés, etc.)
  \item Lenguaje formal (matemática o algún lenguaje de programación o lógico)
  \item Seudocódigo (una mezcla de los dos anteriores)
  \item Como programa de computador
  \item Como hardware (esto es difícil de imaginar en este punto, pero sí)
\end{itemize}
% TODO: Lenguaje formal vs programa de computador, son redundantes?
